\chapter*{Conclusioni}
\addcontentsline{toc}{chapter}{Conclusioni}
In ambito nucleare, le simulazioni numeriche che integrano calcoli neutronici e 
termoidraulici sono di notevole importanza se si considera la complessità delle 
equazioni che governano la fisica del reattore nucleare.
Inoltre le possibilità offerte dalla piattaforma numerica di garantire 
flessibilità e versatilità attraverso la gestione e lo scambio di informazioni tra 
codici differenti consentono di analizzare il comportamento multifisico nel suo 
complesso.

In questa tesi è stato affrontato lo studio dell'accoppiamento neutronico e 
termoidraulico al fine di ottenere una valutazione accurata della distribuzione 
di flusso neutronico, di potenza termica e di temperatura all'interno del 
nocciolo del reattore nucleare LFR.
In particolare, il reattore nucleare preso in considerazione per 
l'implementazione è il reattore dimostrativo ALFRED che si pone lo scopo di 
dimostrare la fattibilità dei reattori veloci raffreddati al piombo; 
quest'ultimi possiedono caratteristiche di sostenibilità e di affidabilità, 
grazie all'utilizzo di uno spettro neutronico veloce e alle proprietà offerte 
dall'uso del piombo liquido come refrigerante.

Al fine di realizzare il modello del reattore, il primo passo è stato la 
creazione del database combustibile contenente le sezioni d'urto 
omogeneizzate sul dominio di un assembly e condensate in 33 gruppi energetici in 
funzione dei principali parametri quali: la temperatura del combustibile, la 
temperatura del refrigerante e il tempo di bruciamento isotopico.
Il database reattoristico è stato ottenuto dalla risoluzione dell'equazione del 
trasporto su ogni assembly che compone il reattore attraverso l'uso del codice 
di cella \textit{Dragon} e delle relative librerie sulle sezioni d'urto 
microscopiche \textit{Jeff3.1.2shem315}.
Il codice di reattore \textit{Donjon} e il codice di termoidraulica 
\textit{FEMuS} cooperano, grazie all'integrazione sulla piattaforma numerica, al 
fine di ottenere un processo iterativo dove avviene lo scambio di informazioni. 
La potenza termica generata nel combustibile è trasferita ai moduli 
di termoidraulica che ne risolvono il campo di temperatura e la nuova 
distribuzione di temperatura è utilizzata per rivalutare il flusso neutronico a 
seguito delle variazioni sulle sezioni d'urto.

Il flusso neutronico è ottenuto dalla risoluzione del sistema di equazioni
$SP_N$ discretizzato a 33 gruppi di energie per ottenere una accurata analisi neutronica. 
Sono stati analizzati i diversi valori della costante moltiplicativa effettiva $k_{eff}$ e 
confrontati con altri studi di ricerca al fine di verificare la validità del 
modello implementato. 
In particolare sono stati valutati in base alle variazioni di temperatura del 
combustibile e all'evoluzione temporale degli isotopi presenti nel combustibile 
a seguito del bruciamento, ottenendo così le variazioni di reattività.
La distribuzione di temperatura sul reattore è stata risolta attraverso 
l'equazione di conservazione dell'energia, in cui si son considerate le diverse 
assembly chiuse e mediate rispetto alla conducibilità termica del piombo.
La temperatura di uscita del refrigerante è risultata accettabile in confronto a quella prevista dal progetto 
ALFRED.

In conclusione i risultati ottenuti sono stati soddisfacenti sia durante le simulazioni di neutronica
sia durante quelle di termoidraulica; in questo contesto l'accoppiamento ha garantito una 
simulazione comprensiva dei principali fenomeni concernenti la fisica del reattore.
Il progetto di tesi può essere sviluppato ulteriormente perfezionando il modello  
termoidraulico che è stato considerato, implementando modelli CFD più sofisticati 
comprendenti anche fenomeni di 
turbolenza per metalli liquidi.
Inoltre si sottolinea che i parametri progettuali reali del progetto ALFRED sono ancora in fase 
preliminare e in continuo cambiamento. In questo contesto, il progetto di tesi
non ha la pretesa di essere una valutazione accurata per il reattore considerato ma si pone piuttosto 
lo scopo di presentare un approccio multifisico computazionale in grado di 
effettuare una simulazione complessiva in ambito reattoristico.
Infatti la possibilità di integrare l'accoppiamento neutronico-termoidraulico 
in un'unica piattaforma numerica può essere un ottimo strumento per la progettazione 
e la valutazione della fattibilità della tecnologia LFR, aprendo prospettive 
future per l'indagine del comportamento multifisico di sistemi energetici complessi. 



